\gddchapter{Refences}

\section{transcript}

Mariusz: W ogóle... O, faktycznie.

Gabriel: Dobra, czekaj, czekaj, czekaj.

Gabriel: Okej, tak, na początku chciałem cię zapytać.

Gabriel: Jak byś się określił jako gracz?

Gabriel: W sensie, jesteś bardziej casual gamer?

Gabriel: Pro gamer?

Gabriel: Coś pomiędzy?

Mariusz: Na ten moment jestem casual gamer.

Mariusz: Nie mam czasu na LV.

Gabriel: Can we crush?

Mariusz: Nie, zawsze, że po prostu jak mam wolny czas,

Mariusz: Spędzę go przed kampem, ale gram w Need for Speed'a ostatnio tylko.

Mariusz: Tak godzinę i tyle.

Mariusz: Dużo nie gram na kampy.

Gabriel: Rozumiem, rozumiem.

Gabriel: I gdy chodzi o asystentów bossowych, to taki nie wiem, Siri albo Aleksa, albo...

Gabriel: Nie używam w ogóle.

Gabriel: Nie?

Gabriel: Dlaczego nie?

Mariusz: Dla mnie są jeszcze niedaleko.

Mariusz: nie na tyle dobrze zrobione, żeby ich pozostać.

Gabriel: Rozumiem.

Gabriel: No moja studia nie potrafi zmieniać postanek, nie?

Mariusz: Także... Moja studia nie rozumie mojego akcentu i źle wybiera numery telefonów.

Gabriel: Ale to już ambitnie.

Gabriel: Ja nigdy nie próbowałem dzwonić przez moją.

Mariusz: Nie, ja próbowałem.

Mariusz: Jakby miałem problem i nie mogłem dać rąk, to chciałem, żeby się przydała, ale...

Mariusz: No.

Mariusz: No i to też głupio, że takie apelowskie, minionowe urządzenie i sztucznie jedynie co nie posiada języka polskiego.

Mariusz: Oj, tak.

Mariusz: Ale bym życzył Divitina.

Mariusz: Co?

Mariusz: Czarny biti ma.

Gabriel: No no ma, nie.

Mariusz: Wiesz, ja premium.

Gabriel: Nie no, po polsku po prostu, no to da znamocieł, wiarnie, to jest polski.

Mariusz: Ale jest taka wersja, że mówisz, jakby, do programu, on pozwoli ci głosem, ale to jest ja premium.

Mariusz: Trzeba kupić.

Gabriel: A, to nie wiem, ja nie używam tego akurat.

Mariusz: Nie, czarny biti miałem akurat w obie wersje.

Mariusz: O.

Mariusz: Jestem zaznajomiony z AI i wiem co ono czuje.

Gabriel: Fantastycznie, to się bardzo przyda.

Gabriel: To powiedzmy jeszcze...

Gabriel: Nie wiem, czekaj.

Gabriel: A tak, czy używasz jakichś takich opcji w stylu... To są takie opcje dostępności na stronach, albo w komputerze, albo w aplikacjach, na przykład czytnik ekranu, albo zmiana kolory, styczki, albo tego typu rzeczy?

Gabriel: Nie?

Mariusz: Nie, nie mam na myśli za bardzo...

Gabriel: No nie wiem, no w LoL-u kiedyś trochę sobie zmieniałem kolory na to, żeby mieć inny schemat dla ludzi, którzy mają daltonizm, żeby lepiej widzieć przeciwników na przykład.

Mariusz: No ale ja nie mam daltonizmu.

Mariusz: No ja też nie, ale...

Mariusz: Ale widzę wszystko chujowo.

Gabriel: No okej, okej, spoko.

Mariusz: Nie, no dobra, kiedyś może...

Mariusz: Jak byłem pro i grałem w Call of Duty to sobie zmieniałem na jaskrawsze kolory, bo podobno kolor turkusowy najbardziej się rzuca w oczy, żeby widzieć bardziej przeciwników.

Mariusz: To się bawiłem.

Mariusz: I zmieniałem ikonki i interfejs.

Mariusz: To było dawno.

Gabriel: Rozumiem, rozumiem.

Gabriel: Ale jakby nie masz takiej potrzeby naturalnej, to robisz to, żeby być lepszy w grze, a nie po to, że inaczej byś nie mógł grać.

Mariusz: Mam, widzę wszystko, to jestem normalny.

Gabriel: Regular.

Gabriel: Regular, tak?

Gabriel: No to powiem ci jeszcze tylko, zanim zaczniemy.

Gabriel: Gra jest grą tekstową, gra jakby...

Gabriel: Idea jest taka, że to jest podtropokalipsa w Warszawie i nie ma zombie.

Gabriel: No wiem, wiem.

Gabriel: I jest tak, że twój brat poszedł w drugiej handlowej pojedzenie dwa dni temu do tej Arkady.

Mariusz: Największe galerii są fajne w Warszawie Największe?

Gabriel: Tu jest Warszawa największa Arkada jest największa?

Gabriel: Arkadia jest największa Nie Arkadia, Arkada Jest Arkada w Warszawie?

Gabriel: Arkada jest pod mostem pojazdowskiego chyba

Mariusz: To jest najbardziej opuszczona galeria handlowa w Warszawie.

Mariusz: Nawet nie wiem, co to jest panorama.

Mariusz: Panorama to jest pierwsza galeria handlowa, jaka powstała po upadku Związku Radzieckiego.

Mariusz: To jest pierwsza galeria, jaka była w Warszawie, a chyba i w Polsce.

Gabriel: Panorama, galeria, poczekaj, co to jest?

Gabriel: Ja nie wiem, w ogóle nigdy tego nie widziałem.

Mariusz: Byłem tam raz, to jest puste.

Mariusz: No, no.

Mariusz: Jakby kurwa dawno to nie byłem, ale...

Mariusz: Nie, ale to jest to.

Gabriel: Ikona luksusu, chłopie.

Mariusz: Tylko, że no na zdjęciach pewnie jest ładniej, ale tam jest pusto.

Mariusz: I sklepy...

Mariusz: No nie masz tam...

Gabriel: No, to wygląda strasznie.

Mariusz: No.

Mariusz: I w ogóle tam jest reklama N coś tam.

Mariusz: Ale to jest sex shop.

Gabriel: No, fajnie.

Mariusz: Nie, nie wiem skąd to jest.

Gabriel: Panorama.

Mariusz: No, ale jednak to jest ślad, w którym stoi, że mamy zachód tutaj u siebie.

Mariusz: I pierwszy panorama, znaczy pierwsza galeria handlowa, więc szkoda, że byśmy tu bywali.

Gabriel: To prawda.

Gabriel: Co jeszcze chciałem powiedzieć?

Gabriel: Jeszcze...

Gabriel: A, dobrze, to jest ta galeria.

Gabriel: No to to jest ta taka inna opuszczona galeria pod mostem.

Gabriel: I co?

Gabriel: I generalnie jeszcze biorąc, to gra to jest...

Gabriel: Chciałem ci pokazać, jak to wygląda.

Gabriel: Możesz tutaj sobie wejść, musisz wejść, musisz wejść.

Ja mogę wszyscy ze sobą wejść.

Mariusz: To jest chyba wodna łóżka.

Gabriel: No, to jest takie wygodne.

Mariusz: A, zakłódka.

Mariusz: Nie chcę iść ku biednemu łóżku.

Gabriel: Dobra, generalnie jedyna różnica względem tego, okej.

Gabriel: Można powiedzieć, tego, to tutaj jest jakaś gra tekstowa.

Gabriel: Tu jest jakby obrazek tła, gdzie jesteś aktualnie.

Gabriel: To jest nazwa lokalizacji, piętro, na którym jesteś.

Gabriel: Tutaj masz opis.

Gabriel: Ten opis jest dosyć długi.

Gabriel: Najważniejsze w opisie jest może pierwsze zdanie i ostatnie zdanie czy dwa zdania.

Gabriel: Jakbyś się czuł zgubiony, to się sugeruj tym.

Gabriel: Resztę tekstu możesz przeskakiwać wzrokiem, nie musisz czytać całości z jakimś wybitnym zrozumieniem.

Gabriel: Co masz w ekipunku?

Gabriel: Nic nie masz aktualnie?

Gabriel: No i to co się różni względu na tego jakby tradycyjnego doświadczenia gamingowego to jest to, że tutaj zamiast sterować za pomocą klawiatury w tej wersji ja bym chciał, żebyś do gry mówił i możesz mówić różnie, możesz mówić używając języka naturalnego, nie musisz mówić jakimiś tam komendami.

Gabriel: I używasz tego kształtowiska R, czyli klikasz R, mówisz to, co chcesz zrobić, później puszczasz i czekasz kilka sekund i wtedy jest taki napis, tu jest aktywowanie, taki bug, za pierwszym razem to się nie wyświetla poprawnie, ale potem już będzie działało.

Gabriel: No i ja sobie będę notował jakieś nowe rzeczy na boku, ale to nie jest to, jak ci dobrze czy źle idzie w tą grę, to w ogóle nie o to chodzi.

Gabriel: Można?

Gabriel: Można wygrać.

Gabriel: O, ok. Zobaczysz jakiś sandbox.

Gabriel: Nie, da się go nawet nie wygrać.

Gabriel: Jakby to jest scenariusz po prostu.

Gabriel: Także...

Gabriel: Teraz dam ci 10 minut po prostu na to.

Gabriel: Przeczytanie tego pizdu?

Gabriel: No, no.

Gabriel: And done.

Gabriel: Nie, no i na grę.

Gabriel: Bo masz 10 minut na całą rozgrywkę.

Gabriel: Mogę ci trochę zwiększyć dubinę tę, bo wiesz, że ta ściąga jest bardzo mała.

Gabriel: Ok. Syk.

Mariusz: Ok, no to jakby to jest bardzo długi opis galerii.

Mariusz: Chciałem ci powiedzieć o tym, żebyś tam był, jak chodzi.

Gabriel: I co mam teraz robić?

Gabriel: No teraz jeżeli chcesz dokonać jakiegoś przykładu gdzieś przejść, tak?

Gabriel: Tu widzisz tutaj, że... Znaczy, nie znam, że są połączone z Small Storage Shed i Clothing Store, tak?

Gabriel: No jeżeli chcesz pójść do tego miejsca, no to klikasz ten przycisk, utrzymujesz go, mówisz, że nie chcesz iść, powiedzmy, i puszczasz, i czekasz chwilę.

Mariusz: Ok. Dm link nie da się.

Mariusz: Tylko tutaj widzę, że nie chcesz mnie naprowadzić, a może mogę robić tak, żeby śledzić te footprinty i zobaczyć skąd, dokąd chodziła ta anomniemana osoba.

Mariusz: Mogę tak zrobić?

Mariusz: Nie możesz?

Gabriel: Mogę oszukać grę?

Gabriel: Nie, możesz bardziej powiedzieć miejsce, do którego idziesz, niestety.

Mariusz: Ok.

Mariusz: A to mam do wyboru jednak coś?

Gabriel: Tak, masz wybór.

Mariusz: Ok, no to pójdzie do tego sklepu z ubraniami, to record i po angielsku.

Gabriel: Tak.

Gabriel: Tylko musisz jakby, teraz poczekaj chwilę, ale musisz kliknąć i tak jak trzymasz na dole, to mówisz.

Mariusz: A, czyli przyciśnięte i mówię.

Gabriel: Tak, tak, tak, dokładnie.

Gabriel: O, i teraz spróbuj jeszcze raz.

Mariusz: I'm going to clothing store.

Gabriel: Teraz jakby przetwarza.

Gabriel: Pierwsze zdanie jest najważniejsze i ostatnie.

Gabriel: 0F to jest jakby piętro, na którym jesteśmy.

Gabriel: Możesz iść w górę i w dół.

Gabriel: To jest zero floor, czyli ground floor, czyli jakby parter w tym wypadku.

Gabriel: Z tego prywatnego miejsca miałeś dwie opcje, w których mogłeś iść, nie?

Gabriel: Mogłeś się cofnąć.

Mariusz: Mogłeś się cofnąć.

Mariusz: I to mi się wyświetli wtedy?

Mariusz: Tak.

Mariusz: Będą obrazki czy napisane?

Gabriel: Napisane.

Mariusz: Wchodzę w cień.

Mariusz: Chciałbym powiedzieć, że...

Mariusz: Yeah, wygrałem.

Mariusz: So I'm going back.

Mariusz: Chciałem się zapytać czy mógłbym z tego miejsca pójść od razu do clothing store.

Gabriel: Niestety nie mógłbyś, niestety w tej wersji gry.

Mariusz: No dobra, to się po prostu wrócę tam i pójdę.

Mariusz: Zacznijmy od czasu.

Mariusz: I'm going to clothing store.

Mariusz: I'm gonna use crowbar on that heavy metal door.

Gabriel: Nie, spokojnie, musisz mi pokazać, widzę wszystko.

Gabriel: Już się napatrzyłem.

Gabriel: Jeszcze raz?

Gabriel: Numerysz się odnowi od Orkster Corridor.

Mariusz: Okej, czyli mogę po prostu teraz mówić, żeby chcę przejść dalej.

Mariusz: Tak, to... Musisz tak, musisz użyć nazwy korytarz.

Mariusz: I'm going to the corridor.

Mariusz: Powinienem tu patrzeć co wygrywam.

Gabriel: Skąd grafiki?

Gabriel: Z AI.

Gabriel: Rysowane lokalnie.

Gabriel: Minęły czas żeby narysować na wszystkich 30 lokalizacji.

Mariusz: I'm going to the bubble tea store.

Mariusz: Nie, ja chcę jechać bubble tea.

Mariusz: Dawno zepsutą i spleśnioną.

Gabriel: Myślę, że one się nie psują.

Mariusz: Nie, ale jakby tam mój brat na pewno by poszedł.

Gabriel: Nie ma go tutaj nadal.

Gabriel: Nie ma go tutaj nadal bubble tea.

Gabriel: Chillować sobie, że tutaj jest bubble tea.

Gabriel: Sorry mordo, dwa dni gdzie nie było.

Gabriel: Zaraz zawoliłem się.

Mariusz: Ok, już jest.

Mariusz: Wracając do food court.

Mariusz: Przynajmniej wiem, że bubble tea ładnie wygląda w środku.

Gabriel: Nie warto, nie warto.

Gabriel: Co tam jeszcze było?

Mariusz: Idę do restauracji.

Gabriel: Ok, teraz, co chciałbym zrobić, to przesadzę Tobie 10 minut.

Gabriel: Zmienimy sterowanie.

Gabriel: Teraz sterujesz za pomocą klautury, a nie za pomocą głosu.

Gabriel: Czyli używasz normalnie klawiszy 1, 2, 3, 4.

Gabriel: Jak masz jakąś akcję, to jest tutaj na dole, np. coś do podniesienia, to klikasz czwórkę.

Gabriel: Więc jakby ten tekst nie masz takiego drugiego znaczenia teraz, no bo masz tutaj wszystkie swoje akcje.

Mariusz: No tak, ale jednak warto zresztą przeczytać.

Mariusz: No jasne, żebyś tam powiedział, co się dzieje.

Mariusz: No widzisz, skąd jestem tylko, raczej ci zbyt długo, ja nawet nie wiem, gdzie je przeczytałem.

Gabriel: Ja, ja.

Mariusz: Prawdopodobnie w ogrodach.

Mariusz: Ale rozumiem, że... Jestem męczachotliwy, więc zaraz sobie... Wpadniesz?

Mariusz: Mogę?

Mariusz: Możesz.

Gabriel: Nieee.

Mariusz: A jak ja będę mieć tam trójkę?

Gabriel: Chciałem zapytać jaką masz strzałeczkę, przy przejściu.

Gabriel: Możesz sobie wybrać jaki chcesz, klikasz ENTER, żeby go użyć.

Gabriel: Już tu nie ma już ludzkiej karty.

Mariusz: Okej, dobrze, dobra.

Mariusz: A crafting?

Gabriel: Crafting, wtedy wybierasz dwa przedmioty ze sobą.

Gabriel: Klikasz tutaj powiedzmy pierwszy i wtedy wybierasz drugi.

Gabriel: I mówi ci, że tego nie można połączyć, czyli crafting polega na łączeniu ze sobą dwóch przedmiotów trzecich.

Gabriel: Okej, okej.

Mariusz: Fajne.

Gabriel: i techniczne, nie wiem czy jesteś na pierwszym piętrze teraz, to jest jakby obrazek, który nie zrobiłem, bo to jest błąd.

Gabriel: Możesz to, co było, tą trawę przejść do innego miejsca.

Gabriel: To może są lounge.

Gabriel: Czy to jest kościół w tle?

Mariusz: Też mi się wydaje, że jakaś msza, ale...

Mariusz: Ale żeby to było słynne?

Mariusz: No jebano, bo daleko chcesz być?

Gabriel: No właśnie, złapu daleko.

Mariusz: Ale tutaj dużo.

Mariusz: No dobra, w kolej.

Mariusz: Tuśmy w małej ziemi.

Mariusz: Stabilizator lega.

Mariusz: Mogi.

Mariusz: Biorę.

Mariusz: Bardzo dobrze.

Mariusz: Żywa, mądra masaż.

Gabriel: No sobie.

Gabriel: To może być drobne jak bubble tea jesteś.

Mariusz: Match this.

Mariusz: Chcesz iść?

Mariusz: Masz... Chcesz?

Gabriel: Prawie nie wzywam na nogi, bo...

Gabriel: Nie wzywam, no ja wiem, no szkoda, wiesz?

Mariusz: Gangster... Okej, tego wcześniej nie było, trzeba mieć stąd...

Mariusz: Chciałem sobie popatrzeć...

Mariusz: W ogóle, kurwa, sklep zbrojny o galerii?

Mariusz: W Polsce?

Mariusz: Dobra, dobra.

Mariusz: Ale jest pistolet Gabrycz, no co mi wkręcać?

Mariusz: Chyba na kulki.

Mariusz: Chyba błąd jest.

Mariusz: Jeszcze raz.

Gabriel: Jest błąd.

Gabriel: Możesz podjeść politycznie ten pistolet.

Gabriel: Szymon podniósł go 8 razy chyba.

Gabriel: Mówię, że zapasowała się pistoletów.

Gabriel: Nie, nie, nie.

Gabriel: Mówię, jednak, jakaś mniejszosta, która odkryła ten błąd, błąd, podniosła pistolet 12 razy i spojrzała mi się prosto w oczy.

Gabriel: No, także... Słuchaj, trzeba być ubezpieczającym.

Gabriel: Czarnobrony nosimy przy sobie 27 pistoletów.

Gabriel: Zawsze.

Mariusz: Czekaj, sprawdzę, czy mogę match tych spłonczyć z empty pistol.

Mariusz: Chyba niestety nie.

Gabriel: Możesz to potem powiedzieć w tym feedbacku pod koniec.

Gabriel: Sobie oszczarowanie przekazać.

Mariusz: The dude.

Mariusz: Dobra, dojdę do inwentory i go dostanę.

Mariusz: Ja ci grę, dobra, dobra.

Gabriel: Dobra, no jakby... Tu jest bug.

Gabriel: Jak klikniesz 4, to nie będzie działać.

Gabriel: Ale jakby uznajemy, że działa, dobra?

Gabriel: Bo tu jest napisane, że może już wstał.

Gabriel: Ale jak teraz pick up your brother, to nie zadziała ten bug.

Gabriel: Tak jak powiedziałem.

Gabriel: No, to jest bug.

Gabriel: Ale teraz musisz z nim wyjść.

Gabriel: Ale mam go, mam go.

Gabriel: Masz go w twoim inwentory, tak?

Gabriel: Musisz z nim wyjść jakoś, ale musisz wyjść...

Gabriel: Sorry.

Gabriel: Inaczej niż wszedłeś do budynku.

Gabriel: To co?

Gabriel: No bo tak ci każę, no.

Mariusz: Ale bo zombie, których nie ma, czy...?

Gabriel: Nie, tu jest bug.

Gabriel: Jakby to się odpaliło, to usłyszałbyś eksplozję i tam by się zapaliło, byś się wejść.

Gabriel: No, a przez to, że jest bug, to nie ma eksplozji i możesz wyjść normalnie, no ale jakby...

Mariusz: No, to co mam robić?

Mariusz: Grać w twoją wizję, czy...?

Gabriel: Tak.

Gabriel: Musisz wyjść inaczej.

Mariusz: Inaczej niż tym razem.

Gabriel: Tak.

Gabriel: Są dwa inne wyjścia z budynku.

Mariusz: Czyli czwórka jest niezostępna teraz.

Mariusz: Te schody.

Gabriel: Nie, schody są dostępne, tylko że jakby samo wejście, to ground entrance, nie?

Gabriel: Tam nie możesz dojść.

Mariusz: Kurwa, ja nie wiem co mi tu...

Gabriel: Tutaj nie wiem czy to jest opisane, ale możesz użyć tego... Czekaj, zobacz, czy to jest w opisie szczegółowości?

Gabriel: A, tak, to tak.

Gabriel: Tu jest np. hint w opisie, co możesz zrobić teoretycznie.

Gabriel: To jest może mało oczywiste, ale tu jest winda i możesz użyć tego keycarda na tej windy.

Gabriel: Tak, no ale już teraz jakby do niej wejść to już też nie... Trochę się rozjechał?

Gabriel: Trochę się rozjechał, aż do teks, ale po trzustku masz elevator.

Gabriel: I elevator jest zbugowany, bo samego elevejtora...

Gabriel: Nieraz się z jakiegoś powodu zjechać niżej.

Gabriel: To jest bug.

Gabriel: Ja tutaj zapauzuję.

Gabriel: Teleportujecie po prostu niżej, dobra?

Mariusz: Dobra.

Gabriel: Bo elewator jedzie jakby... Chyba czy pytanie, czy chcę jechać w górę czy w dół?

Gabriel: Bo elewator jedzie...

Gabriel: No jakby jestem na parterze, więc jechałbym do piwnicy.

Gabriel: Okej.

Gabriel: Dobra, dobra, dobra.

Gabriel: To ja to szybko zrobię tutaj.

Gabriel: Sorry, tak kaszlę.

Gabriel: Mam gardło trochę zasuszone.

Gabriel: Cyk.

Gabriel: Dobra.

Gabriel: No jakby ci... Itemy ci.

Gabriel: No bo restartowaliśmy grę, ale... To nie ma dużego znaczenia.

Mariusz: Nie mam inwentory.

Gabriel: Ale ty jakby...

Gabriel: No jakby teraz...

Gabriel: Ja ci trochę pomogę.

Gabriel: Jeżeli nie masz, to uznajmy, że coś działa, albo nie działa, nie?

Mariusz: Dobra.

Gabriel: Ale no...

Mariusz: Ale mam te rzeczy w domniemaniu.

Gabriel: W domniemaniu masz, tak.

Mariusz: Chcę rozjebać się ułomem samochodu.

Gabriel: No to tylko do błądzi dobrałem.

Mariusz: I pedale.

Mariusz: Znaczy idę do tych samochodów, żeby nie było.

Mariusz: O kurwa, to jakiś Land Rover Mojito miał takiego.

Mariusz: A jak nadzieję, że to nie jest item, to...

Gabriel: Nie ma żadnego item.

Mariusz: No tak, czyli... Tobie mówię, że chciałbym wziąć krobal na samochodzie.

Mariusz: Nie, to się nie stało.

Mariusz: Okej, to nie, to nie.

Mariusz: Tutaj jest portowego.

Mariusz: Okej, tutaj kluczyk.

Mariusz: Też się nie pojawia.

Gabriel: No, jest bug z tym kluczem, nie wiem dlaczego.

Mariusz: No bardzo blisko byłeś generalnie, wiesz?

Gabriel: Tak, tym Jeepem.

Gabriel: Masz ten zamknięty ten?

Gabriel: Jakbyś wziął łom?

Gabriel: To mógłbyś go otworzyć?

Gabriel: Chciałem to zrobić.

Gabriel: No widzisz, no.

Gabriel: Ale nie miałem łomu.

Mariusz: Nie miałeś łomu, bo się zepsuło.

Mariusz: Nie, Bóg mi zabrał górę.

Mariusz: Bóg mnie terpnął na parking podziemny i zabrał mi górę.

Gabriel: I brata mi zabrał.

Gabriel: No słuchaj, nie, ale brat tam jest gdzieś tam w tle.

Gabriel: I byś później mógł odjechać.

Gabriel: Także jakby...

Gabriel: Zakładając, że gra by się nie rozsypała, to byś się skończył 20 minut, nie?

Gabriel: Ok. Dobra.

Gabriel: Chodźcie, teraz zadam kilka pytań.

Mariusz: Ok.

Gabriel: Po polsku.

Gabriel: Ok. O, czekaj.

Mariusz: Czy jestem przesłuchiwany teraz?

Mariusz: Tak, to nasz przesłuchanie.

Mariusz: Dobrze, bo jakś tak ona się obudzi.

Gabriel: No, bardzo dobrze, bo prawda jest.

Gabriel: Dobra, poczekaj.

Gabriel: Dobra, teraz skończyłeś obie wersje, znaczy grałeś obie wersje.

Gabriel: I jak się z tym czujesz?

Gabriel: Znaczy, jak się czujesz?

Mariusz: Okej, yyyy, jak się było?

Mariusz: No jakby druga wersja jest mi bardziej przystępna, bo jestem przyzwyczajony do czegoś takiego, a pierwsza to była taka...

Mariusz: Coś nowego, coś takiego, nic innego.

Mariusz: Chociaż nie, grałem kiedyś w grę z robionym głosem, to był RTS, gdzie wydawałeś komendy głosowe jednostką, może kojarzysz.

Mariusz: Jak się nazywa?

Mariusz: Nie wiem.

Mariusz: W ogóle nie wiem.

Mariusz: Ale gra była fajna.

Gabriel: To stara gra jakaś?

Mariusz: Starsza, powiedzmy, że od okresu 2010 do teraz.

Mariusz: Starsza niż 2010 na pewno nie.

Mariusz: No i była ponadczasowa, bo sterowałeś głosem i wysyłałeś komendę głosową jednostkom, że helikoptery lecą i rzeszterdolają to.

Gabriel: To jest ciekawe, to ja to sobie uwzględniam, to w sumie muszę to znaleźć potem.

Mariusz: Musisz to znaleźć.

Mariusz: To na pewno jest gra psychiczna RTS w czasie rzeczywistym i jednocześnie podłączasz mikrofon, wydajesz im głosowy.

Mariusz: Spoko.

Mariusz: Tylko nie pamiętam czy po polsku, raczej po angielsku.

Mariusz: No poszukam.

Mariusz: Więc kiedyś coś takiego grałem.

Mariusz: Yyyyyyyyyyyyyyyyyyyyyyyyyyyyyyyyyyyyyyyyyyyyyyyyyyyyyyyyyyyyyyyyyyyyyyyyyyyyyyyyyyyyyyyyyyyyyyyyyyyyyyyyyyyyyyyyyyyyyyyyyyyyyyyyyyyyyyyyyyyyyyyyyyyyyyyyyyyyyyyyyyyyyyyyyyyyyyyyyyyyyyyyyyyyyyyyyyyyyyyyyyyyyyyyyyyyyyyyyyyyyyyyyyyyyyyyyyyyyyyyyyyyyyyyyyyyyyyyyyyyyyyyyyyyyyyyyyyyyyyyyyyyyyyyyyyyyyyyyyyyyyyyyyyyyyyyyyyyyyyyyyyyyyyyyyyyyyyyyyyyyyyyyyyyyyyyyyyyyyyyyyyyyyyyyyyyyyyyyyyyyyyyyyyyyyyyyyyyyyyyyyyyyyyyyyyyyyyyyyyyyyyyyyyyyyyyyyyyyyyyyyyyyyy

Mariusz: co możesz zrobić, bo ci daje dostępne opcje, ale jak grasz na gło, to musisz myśleć w tym faktycznie.

Mariusz: Ja też nie lubię pomijać rzeczy, więc... O, żyjesz, żyjesz?

Gabriel: Tak, sorry.

Gabriel: Garbo jest takie zasuszone.

Mariusz: Nie lubię pomijać rzeczy, więc jak widzę dostępne opcje, to będę wiedział, czy mogę do czegoś wrócić, czy mogę coś zrobić, czy czegoś nie zrobiłem.

Gabriel: No, rozumiem, rozumiem, rozumiem.

Mariusz: Więc tak mi się wydaje.

Gabriel: Masz taką większą świadomość możliwości po prostu.

Mariusz: No i to mi się podoba, lubię wiedzieć, kiedy wystaje ta opcja.

Mariusz: To jest takie straszczenie, przez to też, chociaż w pierwszej wersji nie muszę czytać cały tekst, żeby wszystko ogarnąć, w drugiej już się z tym mogę skutknąć, bo widzę gotowe opcje wyborowe, wyboru.

Gabriel: Tak, możesz tutaj też nie skipować tekstu, prawda?

Mariusz: Więc jeśli chodzi o pierwszą wersję, to o wiele lepiej, jeśli chcesz, żeby twój tekst był przeczytany w całości i imersyjny.

Mariusz: Aha.

Mariusz: Ale no, ja jestem zwolennikiem drugiej wersji.

Gabriel: Okej, okej.

Gabriel: A coś cię zaskoczyło, jak grałeś?

Mariusz: Ileś tekstu.

Gabriel: Okej.

Gabriel: Okej, dobra, rozumiem.

Gabriel: Jakbyś miał opisać różnicę pomiędzy jedną wersją, a drugą wersją?

Gabriel: Dla osoby, która nie grała w ani jedną.

Gabriel: To jak byś ją opisał?

Mariusz: Możemy zrobić głosowe odrobienie.

Gabriel: Jak używałeś głosu do nawigacji to czym się kierowałeś przy wyborze tego co powiesz.

Mariusz: Kierowałem się tym jak ja bym wybrał w tej sytuacji.

Gabriel: w których używałeś.

Gabriel: Czyli... to trochę trudne pytanie jest.

Gabriel: Może ja źle formułuję po polsku.

Gabriel: Ale będę myśleć, że na przykład ja mówię do komputera.

Mariusz: A, wiem o co ci chodzi.

Gabriel: W pierwszej osobie.

Gabriel: I dlaczego tak myślę?

Gabriel: Dlaczego tak?

Mariusz: No, jakby używałem najprostszych słów, bo wiem, że to jest AI.

Mariusz: W formie to, że ja to robię, żeby ono to zrozumiało.

Mariusz: Jako, że się utożsamiam z bohaterem, że to jestem ja.

Mariusz: A nie mówię, że he goes there and I'm going there.

Gabriel: Dokładnie o to mi chodziło.

Gabriel: Okej.

Gabriel: I jak to było używać głosu ukrzew?

Mariusz: Dziwnie, to nie jest takie codzienne.

Mariusz: Zwłaszcza przy jakimś...

Gabriel: Siedzi przy łóżku.

Mariusz: Bardzo blisko, więc trochę czuję się oceniający.

Mariusz: Ale to chyba... Tak, rozumiem.

Mariusz: Ja w ogóle jak ktoś obserwuje mnie przy tym co robię, to po prostu...

Mariusz: Dziwnie też gadać do kompa, zawsze, dla mnie.

Mariusz: No, więc takie... ankami.

Gabriel: Okej.

Gabriel: A w której wersji czułeś się bardziej w grze, jakby, w... Tak użyłeś słowa imersyjne, tak.

Gabriel: Co to...

Mariusz: Powiedzmy, że w pierwszej właśnie z głosem akurat, jakby paradoksalnie, no bo musiałem się bardziej skupić i przeczytać wszystko i pomyśleć to, co ja robię i to mówić.

Gabriel: Okej.

Mariusz: Więc powiedzmy tak, że w pierwszej.

Gabriel: Okej.

Gabriel: A jak ten, jakby to powiedzieć po polsku, effort, wysiłek, który musiałeś włożyć w mówienie do komputera, możesz porównać z wysiłkiem, który klikałeś w klawisza.

Mariusz: Jest większy, bo w drugiej wersji mamy gotowe opcje, które można wybrać, a w pierwszej wersji to...

Mariusz: Potrzeba wiedzieć, co chce powiedzieć, żeby nie było tym dużym obłukiem i miało to sens.

Mariusz: Czyli przemyśleć swoje decyzje, a nie z góry wybrać już gotowe przez scenariusz.

Gabriel: Okej.

Gabriel: A czy to, że mogłeś powiedzieć jakby cokolwiek do tego komputera sprawiało, że czułeś się bardziej w kontroli tego, co robisz?

Gabriel: Pod kontrolą tego, co robisz?

Mariusz: Ale nie mogę powiedzieć cokolwiek.

Gabriel: No tak, nie mogłeś, nie?

Gabriel: Ale mogłeś teoretycznie jakby próbować, nie?

Gabriel: A to, że nie miałeś tej listy tych takich rzeczy, tylko mogłeś jakby formułować to, czy to sprawiało, że czułeś się jakbyś miał mniejszą kontrolę, czy sprawiało, że czułeś się bardziej przytłoczony w porównaniu do listy komentów?

Mariusz: Miałem taką, znaczy można odczuć wrażenie, że się ma taką większą kontrolę nad grą, nad tym co się robi, ale ona jest dosyć fałszywa, no bo wiadomo, że nie może zrobić wszystkiego aż tak dobrze jak nie jest zaprojektowane.

Mariusz: Mamy jednak właśnie określone jakieś wybory, które są tam między słowami w opisie.

Mariusz: No ale można uznać, że tak, głosowo czuję się, że mam większe wybory i większą kontrolę.

Gabriel: Jakbyś porównał ogólne doświadczenia przy użyciu tych dwóch wersji?

Gabriel: Odczucia takie, no.

Mariusz: Pytasz mnie już o całą grę?

Gabriel: Tak, holistycznie jakbyś, te dwie wersje jakby się porównał odczuciowo.

Mariusz: Mam wybrać, która lepsza?

Gabriel: Tak, to też twoja droga.

Mariusz: Dla mnie, jako tego już starego biedecza gier i wygodnickiego gościa, to wybrałbym trzecią opcję.

Mariusz: Bardzo mi się dobrze kojarzy, nie wiem czy...

Mariusz: Trzecią?

Mariusz: Drugą.

Mariusz: Pograłem w trzecią, sorry.

Mariusz: Ja nie bałem.

Mariusz: Wybrałbym drugą opcję z tekstem podanym.

Mariusz: Bo mi się tekstowe gry były kiedyś u mnie w serduszku.

Mariusz: Nie wiem, może kojarzysz taką grę flashową przez przeglądarkę o smokach, że się chodzi o swojego smoka.

Mariusz: Nie, nie, nie, nie kojarzę.

Mariusz: To były czasy harcerstwa najwcześniej u przeglądarcy i to była tekstowa gierka.

Mariusz: I wybieraliście właśnie opcję.

Mariusz: Z tym mi się kojarzy, więc podoba mi się to.

Mariusz: Jest to łatwe i wygodne.

Mariusz: Jest.

Mariusz: Jest, to jest moje uczucie.

Mariusz: Okej.

Mariusz: Okej.

Gabriel: I czy jest coś jeszcze na koniec, co chciałbyś dodać, czego jeszcze nie rozmawialiśmy, albo coś od siebie wnieść?

Mariusz: Jakaś krytyka, jakaś opinia?

Mariusz: się nagrywa, bo potem nie bytnie.

Mariusz: Ogólnie fajna, fajna gierka, pomysł.

Mariusz: Lubię takie storyteller.

Mariusz: Mam nadzieję, że będzie więcej romansów.

Mariusz: I jedyna taka uwaga to dużo jest tego tekstu opisującego, może zbytnio co jest dookoła, no bo to jest potrzebne.

Mariusz: Zwłaszcza przy nieznaczących lokacjach, bo to trochę zabiera czasu.

Mariusz: W tym Babułstaurie fajnie.

Mariusz: No chyba, że ma być coś w przyszłości w tym albumie, nie wiem, nie znam swoich planów.

Mariusz: To jeśli coś tam będzie w przyszłości, to okej, ale...

Gabriel: No tak, ale w tej wersji, w której grałeś, to...

Mariusz: W tej wersji, w której grałem, to spoko, fajnie, jakby to buduje świat i tak dalej, ale to nic nie daje, więc... Tak, no jak świat jest...

Gabriel: Jeżeli to nic nie możesz zrobić, no to jakby...

Mariusz: Czasami wolałbym, jeśli przechodzę przez korytarz, który jest tylko łącznikiem między lokacjami, to wolałbym trzy proste zdania, że to jest korytarz i śmierdzi tam gównem.

Mariusz: Tak.

Mariusz: Tyle mi wystarczy.

Mariusz: i więcej konkretów, znaczy więcej konkretów, bardziej takich zarysowanych, to co można wykonać.

Gabriel: Tak, tak, tak, czyli jeżeli coś jest opisane, że dotyka to jakoś mechanizmu w grze, tak?

Mariusz: Tak, nie to, że jest to opis rzeczy, z którą nie można brać pod interakcję, tylko opis rzeczy, którą właśnie mogę wejść w interakcję.

Gabriel: Wszystko co jest opisane jest w jakimś stopniu interaktywne.

Mariusz: Ok.

Mariusz: No i to.

Mariusz: Podoba mi się grafiki.

Mariusz: Przypominają mi te polskie gierki tekstowe.

Mariusz: Tam też były takie lekkie grafiki, żeby to nie było tylko puste.

Mariusz: No i mi się wydaje, że vibe'u, żeby to dopełniało imersji, to przydałoby się udźwiękowanie lekkie do proga.

Mariusz: ale takie dźwięki i ambience, jak jesteś w tej galerii, jakiś szum wiatru między tymi opuszczonymi sklepami, przepięknie kapanie wody, albo też tych podnoszeń rzeczy.

Mariusz: Myślę, że to jest całkiem... dopełniało to fajnie.

Mariusz: Bo tak w ciszę grać, to...

Mariusz: Pełna zgoda.

Gabriel: A miejsce byłoby fajne, to prawda.

Gabriel: No dobra, dziękuję ci bardzo.



